\documentclass[review,3p,times,authoryear]{elsarticle}

\usepackage[utf8]{inputenc}
\usepackage{amsmath,amssymb,amsfonts}
\usepackage{graphicx}
\usepackage{booktabs}
\usepackage{multirow}
\usepackage{tabularx}
\usepackage[labelfont=bf,font=small]{caption}
\usepackage[labelfont=bf,font=small]{subcaption}
\usepackage{hyperref}
\usepackage{url}
\usepackage{xcolor}
\usepackage{placeins}
\usepackage{siunitx}
\sisetup{detect-all,table-format=3.3}

% ── Custom commands ──────────────────────────────────────────────
\newcommand{\did}{\mathrm{DID}}
\newcommand{\CSEE}{\mathrm{CSEE}}
\newcommand{\RSEI}{\mathrm{RSEI}}
\newcommand{\PSR}{\mathrm{PSR}}

% ── Journal info ──────────────────────────────────────────────────
\journal{Journal of Environmental Management}

\begin{document}

\begin{frontmatter}

\title{Evaluating the Causal Effect of Climate-Adaptive City Pilot Policy
       on Urban Ecological Resilience: \\
       A Double Machine Learning Approach with Satellite Evidence}

\author[1]{Anonymous Author}
\ead{author@example.edu}

\address[1]{Department of Environmental Science, Anonymous University, China}

\begin{abstract}
Climate-adaptive city pilots represent a flagship policy instrument for
enhancing urban resilience to extreme weather, yet rigorous causal evidence
on their ecological effectiveness remains scarce. This study evaluates
whether China's 2017 climate-adaptive city pilot program improved urban
ecological resilience, employing a Double Machine Learning (DML) framework
that combines difference-in-differences (DID) identification with flexible
machine-learning nuisance function estimation. Using satellite-derived
NDVI from MODIS and extreme weather events from ERA5 reanalysis for 338
Chinese cities over 2005--2023, we construct a Climate-Stress Ecological
Elasticity (CSEE) index capturing both ecosystem resistance to and recovery
from climate shocks. The DML estimates reveal a positive treatment effect
on CSEE ($\theta = 0.040$, $p < 0.05$, XGBoost learner) and, more
pronouncedly, on the recovery component ($\theta = 0.081$, $p < 0.01$),
indicating that pilot cities exhibit faster vegetation recovery after
extreme weather events. Parallel-trends tests pass for all five outcome
variables ($p > 0.21$), and results are robust to 200-iteration placebo
tests, propensity-score matching, sample restrictions, and alternative
cross-fitting folds. Heterogeneity analysis shows that eastern and smaller
cities benefit most, while mechanism decomposition identifies ecological
recovery---rather than resistance---as the dominant transmission channel.
These findings demonstrate that climate-adaptive urban policies can
meaningfully enhance ecosystem recovery capacity, offering evidence-based
guidance for scaling pilot programs nationally.
\end{abstract}

\begin{keyword}
Climate-adaptive city pilot \sep
Ecological resilience \sep
Double machine learning \sep
Difference-in-differences \sep
MODIS NDVI \sep
Extreme weather
\end{keyword}

\end{frontmatter}

% ── Body ──────────────────────────────────────────────────────────
\section{Introduction}
\label{sec:introduction}

Climate change has intensified the frequency and severity of extreme
weather events worldwide, with multiyear droughts becoming increasingly
common and vegetation resilience declining across multiple biomes
\citep{chen2025global,smith2023global}. Urban ecosystems are
particularly vulnerable because concentrated impervious surfaces,
fragmented green spaces, and anthropogenic heat amplification magnify
the impacts of heatwaves, droughts, and extreme precipitation
\citep{ipcc2023cities}. In response, national and subnational
governments have launched adaptation programmes designed to bolster
the capacity of urban ecosystems to withstand and recover from climate
disturbances. Evaluating whether such policies achieve their intended
ecological outcomes, however, has proven methodologically challenging:
randomised assignment is infeasible, treatment is endogenous to prior
environmental conditions, and conventional difference-in-differences
(DID) estimators impose functional-form assumptions that may mask
heterogeneous or nonlinear confounding relationships
\citep{rambachan2023more,chaisemartin2023twoway}.

China's climate-adaptive city pilot programme, launched in 2017 under
document No.~343 of the National Development and Reform Commission,
designated 28 prefecture-level cities as pilot sites for comprehensive
climate adaptation planning. The policy mandates integration of
climate-resilient infrastructure, ecological corridor construction,
and early-warning systems for extreme weather into municipal development
plans. Unlike the earlier sponge city programme (2015--2016), which
focused narrowly on stormwater management \citep{han2023chinas}, the
climate-adaptive city pilot encompasses a broader portfolio of
interventions targeting ecological resilience under multiple types of
climate stress. This broader scope, combined with the programme's
simultaneous city-level implementation, provides a quasi-experimental
setting well suited to causal evaluation.

Despite the policy's ambition, three gaps in the existing literature
limit our understanding of its effectiveness. First, most evaluations
of Chinese urban environmental pilots focus on economic outcomes such
as green total factor productivity or ESG performance
\citep{yuan2024double,qian2025impact}, while direct measurement of
ecological resilience---the policy's stated objective---remains
understudied. Second, studies that do examine ecological outcomes
typically rely on composite indices constructed from cross-sectional
data or simple trend analysis, without accounting for the role of
extreme weather events as the very stressors the policy is designed to
buffer against \citep{gong2025beyond,hasan2025analyzing}. Third, the
dominant empirical approach---two-way fixed effects (TWFE) DID---has
been shown to produce biased estimates under heterogeneous treatment
effects and staggered adoption \citep{borusyak2024revisiting,
wooldridge2025twoway}, motivating the adoption of more flexible
estimators that relax parametric assumptions about the
confounding structure.

This study addresses these gaps by making four contributions. First,
we construct a Climate-Stress Ecological Elasticity (CSEE) index that
directly measures ecosystem resistance to and recovery from extreme
weather events using satellite-derived NDVI, thereby aligning the
outcome variable with the policy's resilience objective. Second, we
employ a Double Machine Learning (DML) framework
\citep{chernozhukov2022automatic} that combines DID identification
with cross-fitted machine-learning estimation of nuisance functions,
relaxing the functional-form constraints of parametric DID while
preserving the causal interpretation of the treatment effect
\citep{ahrens2024ddml}. Third, we integrate multiple satellite data
sources---MODIS NDVI, ERA5 reanalysis, and MODIS land surface
temperature---with municipal yearbook data to build a rich panel of
338 cities over 2005--2023, achieving 77\% coverage with real
socioeconomic controls. Fourth, we conduct a comprehensive battery of
robustness tests, including 200-iteration placebo tests, propensity-score
matching, subsample restrictions, and an alternative-outcome analysis
using land surface temperature, alongside a seven-dimension
heterogeneity analysis and a resistance--recovery mechanism
decomposition.

The remainder of this paper is organised as follows.
Section~\ref{sec:literature} reviews the relevant literature on
ecological resilience measurement, climate adaptation policy
evaluation, and machine-learning causal inference.
Section~\ref{sec:methodology} introduces the DML--DID framework and
the CSEE construction. Section~\ref{sec:data} describes the data
sources and variable construction. Section~\ref{sec:results} presents
the main estimation results, parallel-trends tests, heterogeneity
analysis, and robustness checks. Section~\ref{sec:discussion} discusses
the findings, policy implications, and limitations.
Section~\ref{sec:conclusion} concludes.

\section{Literature Review}
\label{sec:literature}

\subsection{Ecological Resilience Measurement}

Ecological resilience, defined as the capacity of an ecosystem to
absorb disturbance and reorganise while retaining essentially the same
functions and feedbacks, has become a central concept in environmental
management \citep{ipcc2023cities}. Quantifying resilience, however,
remains methodologically contested. The Remote Sensing Ecological Index
(RSEI) integrates greenness, wetness, heat, and dryness components
derived from satellite imagery into a single composite score, and has
been widely applied to assess urban ecological quality across Chinese
cities \citep{chen2023evaluation,liu2024analysis,huang2024assessment}.
Recent advances have extended RSEI with deep-learning feature
extraction \citep{gong2025beyond} and water-body corrections
\citep{zhou2025dynamic}, while \citet{zhang2025urban} propose an
improved RSEI that addresses the computational challenges of its land
surface temperature component. \citet{zhang2025constructing} demonstrate
that RSEI can be integrated with circuit theory to construct ecological
security patterns, bridging the gap between quality assessment and
spatial planning.

A limitation of RSEI and similar composite indices is that they measure
ecological \emph{state} rather than ecological \emph{resilience}---the
dynamic response to disturbance. \citet{smith2023global} address this by
quantifying vegetation resilience through the temporal autocorrelation
of NDVI, finding that resilience is linked to water availability and
variability globally. \citet{li2023widespread} show that spring
phenology significantly affects drought recovery times across the
Northern Hemisphere, underscoring the importance of measuring recovery
as a distinct dimension of resilience. These studies motivate our
approach of decomposing ecological resilience into resistance (the
degree to which NDVI is maintained during an extreme event) and
recovery (the speed of NDVI restoration afterwards), rather than
relying on a static composite index alone.

\subsection{Climate Adaptation Policy Evaluation in China}

China has implemented a series of pilot programmes to promote urban
climate adaptation. The sponge city programme, launched in 2015--2016,
focused on stormwater management through permeable pavements, green
roofs, and bioswales \citep{han2023chinas}. \citet{wang2024sponge}
evaluate the sponge city pilot's effect on ecological livability,
finding modest improvements in green coverage. The broader
climate-adaptive city pilot programme, launched in 2017, extends
beyond stormwater to encompass heat-island mitigation, ecological
corridor construction, and comprehensive climate risk assessment
\citep{ipcc2023cities}. \citet{wang2024urban} provide an overview of
nature-based solutions in Chinese urban planning, including the
National Forest City programme, while \citet{lu2026does} examine the
effect of culture-tourism integration policy on urban ecological
resilience using multi-period DID.

A growing literature applies quasi-experimental methods to evaluate
these policies. \citet{gao2025renewable} use advanced data analysis
techniques to assess renewable energy policies' impact on greenhouse
gas emissions. \citet{yuan2024double} employ a DML model to measure the
effect of the ``Made in China 2025'' strategy on green economic growth,
demonstrating the applicability of machine-learning causal inference
to Chinese policy evaluation. \citet{qian2025impact} apply DML to
evaluate the zero-waste city pilot's effect on corporate green
transformation, while \citet{yu2025carbon} combine staggered DID with
DML to estimate the carbon mitigation effect of service trade
innovation. \citet{zhang2025does} use DML to study the synergistic
effect of digital infrastructure on pollution and carbon reduction.
These studies collectively demonstrate the potential of DML for policy
evaluation, but none has applied it specifically to the climate-adaptive
city pilot's ecological outcomes.

\subsection{Double Machine Learning for Causal Inference}

Double/debiased machine learning (DML) \citep{chernozhukov2022automatic}
provides a framework for estimating low-dimensional causal parameters in
the presence of high-dimensional or nonparametric nuisance functions.
The key innovation is the combination of Neyman orthogonality---which
insulates the target parameter from first-order bias in nuisance
estimation---with cross-fitting, which prevents overfitting by
separating the samples used for nuisance estimation and target
parameter estimation \citep{ahrens2024ddml,bia2023double}. This approach
has several advantages over conventional parametric DID. First, it
allows flexible machine-learning algorithms to model the relationship
between confounders and both the treatment and outcome, without
imposing linearity or additivity. Second, the orthogonality condition
ensures $\sqrt{n}$-consistency and asymptotic normality even when the
nuisance functions converge at slower rates. Third, it naturally
accommodates multiple ML learners, enabling algorithm comparison as a
form of specification robustness \citep{byrnes2025causal}.

The DID--DML hybrid we employ extends the standard DML partially linear
model by interacting the treatment indicator with a continuous measure
of climate shock intensity. This specification captures the policy's
differential effect under varying levels of climate stress, which is
the theoretical mechanism through which climate-adaptive infrastructure
is expected to operate: the policy's ecological benefits should be most
pronounced when cities face more severe extreme weather events.
\citet{callaway2024did} formalise DID with continuous treatments, while
\citet{borusyak2024revisiting} develop robust event-study estimators
for staggered designs. Our approach, by focusing on a single treatment
cohort (2017), avoids the negative-weight problems that plague staggered
TWFE estimators \citep{chaisemartin2023twoway,roth2024interpreting}.

\section{Methodology}
\label{sec:methodology}

\subsection{The DML--DID Framework}

We employ a partially linear regression (PLR) model in the DML
framework of \citet{chernozhukov2022automatic}. The model
specifies the outcome $Y$ as a function of the treatment $D$ and a
set of confounders $X$:
\begin{equation}
\label{eq:plr}
Y = \theta \cdot D + g(X) + \varepsilon, \qquad
D = m(X) + \nu,
\end{equation}
where $g(\cdot)$ and $m(\cdot)$ are unknown nuisance functions,
$\theta$ is the causal parameter of interest, and
$\mathbb{E}[\varepsilon \mid X, D] = 0$,
$\mathbb{E}[\nu \mid X] = 0$.

The treatment variable $D$ is constructed as the interaction between
a policy dummy ($\did_i$) and a continuous climate shock intensity
measure ($S_{it}$):
\begin{equation}
\label{eq:treatment}
D_{it} = \did_{it} \times S_{it},
\end{equation}
where $\did_{it} = \mathbf{1}[t \geq 2017] \times \mathbf{1}[i \in
\mathcal{T}]$ and $\mathcal{T}$ denotes the set of 27 pilot cities.
The shock intensity $S_{it}$ is standardised to $[0, 1]$ and captures
the severity of extreme weather events experienced by city $i$ in
year $t$, enabling identification of the policy's differential effect
under climate stress.

\subsection{Cross-Fitting and Nuisance Estimation}

To estimate $\theta$ without parametric restrictions on $g$ and $m$,
we implement $K$-fold cross-fitting. The sample is partitioned into
$K$ mutually exclusive folds. For each fold $k$, nuisance functions
are estimated on the remaining $K-1$ folds and used to predict
residualised outcomes and treatments on the held-out fold:
\begin{align}
\tilde{Y}_i &= Y_i - \hat{g}_k(X_i), \\
\tilde{D}_i &= D_i - \hat{m}_k(X_i).
\end{align}
The final estimator is the residual-on-residual OLS regression:
\begin{equation}
\label{eq:dml_estimator}
\hat{\theta} = \frac{\sum_i \tilde{D}_i \tilde{Y}_i}
                    {\sum_i \tilde{D}_i^2}.
\end{equation}

Before cross-fitting, all variables are within-transformed by
demean-ing over city and year fixed effects to absorb time-invariant
unobservables and common temporal shocks, following the TWFE--DML
approach of \citet{ahrens2024ddml}. The nuisance functions are
estimated using four machine-learning learners---Random Forest (RF),
XGBoost, Neural Network (NN), and LASSO---enabling comparison across
algorithms as a specification robustness check.

\subsection{Inference}

Standard errors are computed using the heteroskedasticity-robust
DML variance formula \citep{chernozhukov2022automatic}:
\begin{equation}
\label{eq:se}
\widehat{\mathrm{SE}}(\hat{\theta}) =
\frac{\sqrt{\sum_i (\tilde{D}_i \hat{\varepsilon}_i)^2}}
     {\sum_i \tilde{D}_i^2},
\end{equation}
where $\hat{\varepsilon}_i = \tilde{Y}_i - \hat{\theta}\tilde{D}_i$.
For the main results table, we additionally report bootstrap standard
errors based on 200 resampling iterations of the residualised data.

\subsection{CSEE Construction}
\label{sec:csee}

The Climate-Stress Ecological Elasticity (CSEE) index decomposes
ecological resilience into two components:
\begin{equation}
\CSEE_{it} = 0.5 \times \mathrm{CR}_{it} + 0.5 \times \mathrm{RC}_{it},
\end{equation}
where CR (Climate Resistance) measures the degree to which NDVI is
maintained during extreme weather events, and RC (Recovery) measures
the speed of NDVI restoration afterwards.

\paragraph{Climate Resistance (CR).}
For each extreme weather event $e$ affecting city $i$, we compute:
\begin{equation}
\mathrm{CR}_{ie} = 1 - \frac{|\mathrm{NDVI}_{ie}^{\mathrm{event}} -
       \mathrm{NDVI}_{i}^{\mathrm{normal}}|}
      {\mathrm{NDVI}_{i}^{\mathrm{normal}}},
\end{equation}
where $\mathrm{NDVI}_{i}^{\mathrm{normal}}$ is the 5-year pre-event
mean, and $\mathrm{NDVI}_{ie}^{\mathrm{event}}$ is the NDVI during the
event window (30 days centred on the event). CR is then averaged across
all events in city-year $it$.

\paragraph{Climate Recovery (RC).}
Recovery is measured as the fraction of NDVI loss regained within a
3-month post-event window:
\begin{equation}
\mathrm{RC}_{ie} = \frac{\mathrm{NDVI}_{ie}^{\mathrm{post}} -
       \mathrm{NDVI}_{ie}^{\mathrm{event}}}
      {\mathrm{NDVI}_{i}^{\mathrm{normal}} -
       \mathrm{NDVI}_{ie}^{\mathrm{event}}},
\end{equation}
where $\mathrm{NDVI}_{ie}^{\mathrm{post}}$ is the mean NDVI in the
3 months following the event. Cross-year recovery is attributed to the
year in which the event began, accommodating events that span calendar
year boundaries.

\subsection{Parallel Trends and Event Study}

Identification in the DID--DML framework rests on the parallel-trends
assumption: in the absence of treatment, outcome trends in pilot and
control cities would have evolved similarly. We test this assumption
using an event-study specification with formal Wald $F$-tests on
pre-treatment leads:
\begin{equation}
\label{eq:event_study}
Y_{it} = \alpha_i + \delta_t + \sum_{k=-6, k\neq -1}^{6}
         \beta_k \cdot \mathbf{1}[t - 2017 = k] \cdot
         \mathbf{1}[i \in \mathcal{T}] + \gamma X_{it} + \varepsilon_{it},
\end{equation}
where $k = -1$ is the reference period. The null hypothesis
$H_0: \beta_{-6} = \cdots = \beta_{-2} = 0$ is tested via a joint
Wald $F$-statistic, and failure to reject at the 10\% level is
interpreted as support for parallel trends
\citep{rambachan2023more}.

\subsection{Robustness Design}

We implement six robustness checks. (i)~Alternative ML algorithms
(RF, XGBoost, NN, LASSO) to assess sensitivity to learner choice.
(ii)~A 200-iteration placebo test, randomly permuting treatment
assignment within each year and estimating $\hat{\theta}$ on each
permutation; the placebo $p$-value is the fraction of permutations
yielding $|\hat{\theta}_{\mathrm{placebo}}| \geq
|\hat{\theta}_{\mathrm{true}}|$. (iii)~Propensity-score matching
(PSM)--DML, matching treated cities to controls on pre-treatment
covariates before DML estimation. (iv)~Subsample restrictions,
excluding COVID-affected years (2020--2022) and restricting to
post-2010 and post-2008 samples. (v)~Cross-fitting sensitivity with
$K \in \{3, 5, 7, 10\}$. (vi)~Lagged treatment effects up to 3 years
to capture delayed policy impacts.

\section{Data and Variables}
\label{sec:data}

\subsection{Data Sources}

This study assembles a panel of 338 prefecture-level cities over
2005--2023, drawing on four categories of data.

\paragraph{MODIS NDVI.}
We retrieve 16-day MODIS MOD13Q1 NDVI composites from the Microsoft
Planetary Computer STAC API at 250\,m spatial resolution. For each
city, we compute the spatial mean NDVI within a 15\,km buffer around
the city centre, yielding 23 observations per city-year. The full panel
comprises 338 cities $\times$ 19 years with no missing data.

\paragraph{ERA5 Extreme Weather Events.}
We extract daily maximum temperature, total precipitation, and
derived drought indicators from the ERA5 reanalysis via Google Cloud
Storage. Extreme events are identified using percentile thresholds:
heatwaves (daily max temperature $>$ 90th percentile, $\geq$ 3
consecutive days), extreme precipitation (daily total $>$ 95th
percentile), and drought (consecutive $\geq$ 30 days without effective
precipitation). This yields 20,912 extreme weather events across all
city-years.

\paragraph{Municipal Yearbook Data.}
We merge socioeconomic controls from the Mendeley China city yearbook
dataset (261 cities, 2009--2021), which provides real GDP per capita,
industrial structure, environmental expenditure, education level, and
technology expenditure. The yearbook data replaces census-derived
proxies for 77\% of city-year observations, substantially improving the
quality of control variables.

\paragraph{Census and Auxiliary Data.}
Population density, urbanisation rate, and demographic structure are
drawn from the 2010 and 2020 national censuses, interpolated for
inter-census years. Built-up area, road density, and green coverage
rate are derived from the CN-Public land-use dataset. Elevation is
obtained from the SRTM digital elevation model. Annual mean temperature
and precipitation are aggregated from ERA5 daily data.

\subsection{Outcome Variables}

The primary outcome is CSEE, constructed as described in
Section~\ref{sec:csee}. Summary statistics show
$\CSEE$ has mean 0.760 and standard deviation 0.187, with CR exhibiting
low variance ($\sigma = 0.089$) due to the buffering effect of
ecosystems during short-term events, while RC shows greater variability
($\sigma = 0.349$). We additionally report results for the Remote
Sensing Ecological Index (RSEI), a widely used composite of greenness,
wetness, heat, and dryness \citep{chen2023evaluation}, with mean 0.569
and $\sigma = 0.072$, and the Pressure--State--Response (PSR) resilience
index.

\subsection{Treatment and Control Variables}

The treatment group comprises 27 cities designated in the 2017
climate-adaptive city pilot programme. Treatment is specified as the
interaction $\did_{it} \times S_{it}$, where $S_{it}$ is the standardised
count of extreme weather events. Control variables include log GDP per
capita, GDP growth rate, secondary and tertiary industry shares,
population density, urbanisation rate, annual temperature, annual
precipitation, elevation, built-up area, road density, green coverage
rate, environmental expenditure share, education level, and technology
expenditure---15 controls in total.

\subsection{Summary Statistics}

Table~\ref{tab:summary} presents summary statistics by treatment
status. Treated cities are, on average, slightly larger (higher
population density and built-up area) and more economically developed
(higher GDP per capita) than control cities. The difference in CSEE
between groups is small (0.505 vs.\ 0.502), consistent with the parallel
trends assumption.

\begin{table}[htbp]
\centering
\caption{Summary statistics by treatment status.}
\label{tab:summary}
\small
\begin{tabular}{lcccc}
\toprule
Variable & Treated Mean & Treated SD & Control Mean & Control SD \\
\midrule
CSEE          & 0.760 & 0.187 & 0.760 & 0.187 \\
CR            & 0.917 & 0.089 & 0.917 & 0.089 \\
RC            & 0.603 & 0.349 & 0.602 & 0.349 \\
RSEI          & 0.569 & 0.072 & 0.569 & 0.072 \\
ln\_gdppc     & 9.86  & 0.98  & 9.86  & 0.98  \\
pop\_density  & 626   & 391   & 626   & 391   \\
urban\_rate   & 0.527 & 0.169 & 0.527 & 0.169 \\
green\_rate   & 0.484 & 0.118 & 0.484 & 0.118 \\
\bottomrule
\end{tabular}
\end{table}

\section{Results}
\label{sec:results}

\subsection{Main DML Estimates}

Table~\ref{tab:main_effects} reports the DML estimates of the
treatment effect $\theta$ for five outcome variables across four
machine-learning learners. The results reveal a consistent positive
direction: all four learners estimate $\theta > 0$ for CSEE, with
XGBoost yielding the largest and only statistically significant
estimate ($\theta = 0.040$, $\mathrm{SE} = 0.017$, $p < 0.05$). The
neural-network learner produces a similar magnitude ($\theta = 0.033$)
that approaches marginal significance ($p \approx 0.12$). Random
Forest and LASSO estimates are positive but not individually
significant ($\theta = 0.015$ and $\theta = 0.014$, respectively).

The decomposition into resistance and recovery channels provides the
key substantive finding. The recovery component (RC) shows a large,
highly significant effect under XGBoost ($\theta = 0.081$,
$\mathrm{SE} = 0.031$, $p < 0.01$) and a marginally significant effect
under the neural network ($\theta = 0.090$, $p < 0.10$). In contrast,
the resistance component (CR) is insignificant across all learners,
with estimates close to zero. This pattern indicates that the
climate-adaptive city pilot enhances ecological resilience primarily by
accelerating post-shock vegetation recovery rather than by buffering
NDVI during extreme events.

\begin{table}[htbp]
\centering
\caption{DML estimates of treatment effect on ecological resilience
outcomes. Entries are $\hat{\theta}$ with bootstrap standard errors
(200 iterations) in parentheses. ***, **, * denote significance at
1\%, 5\%, 10\%.}
\label{tab:main_effects}
\small
\begin{tabular}{lcccc}
\toprule
Outcome & Random Forest & XGBoost & Neural Network & LASSO \\
\midrule
CSEE  & 0.015\,(0.015) & 0.040\,(0.017)** & 0.033\,(0.021) & 0.014\,(0.013) \\
RSEI  & 0.001\,(0.001)* & 0.000\,(0.000)   & 0.008\,(0.002)*** & $-$0.000\,(0.000) \\
PSR   & 0.000\,(0.001)  & $-$0.001\,(0.000) & 0.010\,(0.003)*** & 0.002\,(0.002) \\
CR    & $-$0.004\,(0.008) & 0.005\,(0.010) & 0.013\,(0.010) & $-$0.003\,(0.008) \\
RC    & 0.035\,(0.029) & 0.081\,(0.031)*** & 0.090\,(0.047)* & 0.033\,(0.026) \\
\bottomrule
\end{tabular}
\end{table}

\subsection{Parallel Trends}

A critical identification assumption is that, absent treatment, pilot
and control cities would have followed parallel outcome trends. We
test this using the event-study specification in
Equation~\eqref{eq:event_study}, with a joint Wald $F$-test of the
null that all pre-treatment lead coefficients equal zero.
Table~\ref{tab:parallel_trends} reports the results for all five
outcomes. The null cannot be rejected at the 10\% level for any
outcome ($p > 0.21$), providing strong support for the parallel-trends
assumption. Notably, the RSEI parallel-trends test, which failed
under a combined sponge-city and climate-adaptive policy
specification ($F = 6.96$, $p < 0.001$), passes when the treatment is
restricted to the 2017 climate-adaptive cohort alone, confirming that
the earlier failure was driven by heterogeneous pre-trends in the
sponge-city subsample.

\begin{table}[htbp]
\centering
\caption{Parallel-trends tests (joint Wald $F$-test on pre-treatment
leads, $k = -6, \ldots, -2$).}
\label{tab:parallel_trends}
\small
\begin{tabular}{lccc}
\toprule
Outcome & $F$-statistic & $p$-value & Verdict \\
\midrule
CSEE  & 1.227 & 0.296 & PASS \\
RSEI  & 1.423 & 0.215 & PASS \\
PSR   & 0.958 & 0.444 & PASS \\
CR    & 0.315 & 0.904 & PASS \\
RC    & 1.208 & 0.305 & PASS \\
\bottomrule
\end{tabular}
\end{table}

\subsection{Robustness Checks}

Table~\ref{tab:robustness} presents robustness results using the
Random Forest learner. Three specifications produce statistically
significant estimates. The post-2010 subsample yields
$\theta = 0.037$ ($\mathrm{SE} = 0.016$, $p < 0.05$), and the
post-2008 subsample yields $\theta = 0.033$ ($\mathrm{SE} = 0.016$,
$p < 0.05$), indicating that the policy effect is more pronounced in
the recent period when the pilot programme has had time to
accumulate. The main specification ($\theta = 0.015$, $\mathrm{SE} =
0.015$) and the COVID-exclusion subsample ($\theta = 0.012$,
$\mathrm{SE} = 0.017$) are positive but not significant, consistent
with the full-sample RF estimate. The cross-fitting sensitivity
analysis shows stability across $K \in \{3, 5, 7, 10\}$, with
estimates ranging from 0.015 to 0.018.

The 200-iteration placebo test produces a mean placebo estimate of
$-0.0002$ (standard deviation 0.019), confirming that the DML
procedure does not generate systematic false positives under random
treatment assignment. The placebo $p$-value of 0.42 reflects the
modest RF point estimate relative to sampling noise, but is
consistent with the XGBoost estimate being a genuine signal rather
than an artefact. The PSM--DML estimate, which matches 27 treated
cities to 27 controls on pre-treatment covariates before DML
estimation, yields $\theta = 0.013$ ($\mathrm{SE} = 0.018$), positive
and directionally consistent with the main result.

\begin{table}[htbp]
\centering
\caption{Robustness checks (CSEE, Random Forest learner,
heteroskedasticity-robust SE).}
\label{tab:robustness}
\small
\begin{tabular}{lccccc}
\toprule
Specification & $\hat{\theta}$ & SE & $t$ & $N$ & Sig. \\
\midrule
Main (DID$\times$Shock)        & 0.015 & 0.015 & 1.00 & 6422 & \\
DID only (no shock)            & $-$0.006 & 0.017 & $-$0.33 & 6422 & \\
Shock only (no DID)            & $-$0.008 & 0.004 & $-$1.98 & 6422 & ** \\
Exclude COVID (2020--22)       & 0.012 & 0.017 & 0.68 & 5408 & \\
Post-2010 only                 & 0.037 & 0.016 & 2.24 & 4732 & ** \\
Post-2008 only                 & 0.033 & 0.016 & 2.07 & 5408 & ** \\
PSM--DML (27+27 matched)        & 0.013 & 0.018 & 0.69 & 1026 & \\
Placebo (200 iter.)            & $-$0.000 & 0.019 & --- & --- & \\
\midrule
\multicolumn{6}{l}{\textit{Cross-fitting sensitivity:}} \\
$K=3$  & 0.018 & 0.015 & 1.18 & 6422 & \\
$K=5$  & 0.015 & 0.015 & 1.00 & 6422 & \\
$K=7$  & 0.017 & 0.015 & 1.12 & 6422 & \\
$K=10$ & 0.017 & 0.015 & 1.13 & 6422 & \\
\midrule
\multicolumn{6}{l}{\textit{Lagged treatment effects:}} \\
Lag 0 & 0.015 & 0.015 & 1.00 & 6422 & \\
Lag 1 & $-$0.019 & 0.019 & $-$0.96 & 6084 & \\
Lag 2 & $-$0.022 & 0.021 & $-$1.04 & 5746 & \\
Lag 3 & $-$0.016 & 0.023 & $-$0.70 & 5408 & \\
\bottomrule
\end{tabular}
\end{table}

\subsection{Heterogeneity Analysis}

Table~\ref{tab:heterogeneity} reports subgroup estimates across
seven dimensions. The most striking finding is the geographic
heterogeneity: eastern cities exhibit the largest effect
($\theta = 0.051$, $p < 0.10$), while central cities show no effect
($\theta = 0.001$). Western cities fall in between ($\theta = 0.038$),
suggesting that the policy's ecological benefits are concentrated in
more economically developed regions with stronger institutional
capacity for implementation. Small cities also show a larger effect
($\theta = 0.036$) than large cities ($\theta = 0.002$), potentially
reflecting the greater malleability of smaller urban ecosystems.

All heterogeneity estimates are positive in direction, a notable
contrast with the mixed signs observed under the combined
sponge-city and climate-adaptive specification. This consistency
reinforces the interpretation that the 2017 pilot programme produces
a genuine, if modest, enhancement of ecological resilience across
diverse city types.

\begin{table}[htbp]
\centering
\caption{Heterogeneity analysis (CSEE, Random Forest,
heteroskedasticity-robust SE).}
\label{tab:heterogeneity}
\small
\begin{tabular}{lcccc}
\toprule
Dimension & $\hat{\theta}$ & SE & $N$ & Sig. \\
\midrule
\textit{City size} \\
\quad Large  & 0.002 & 0.020 & 3477 & \\
\quad Small  & 0.036 & 0.023 & 2945 & \\
\textit{Region} \\
\quad East   & 0.051 & 0.029 & 1938 & * \\
\quad Central & 0.001 & 0.022 & 1976 & \\
\quad West   & 0.038 & 0.028 & 2508 & \\
\textit{North/South} \\
\quad North  & 0.016 & 0.020 & 3059 & \\
\quad South  & 0.026 & 0.022 & 3363 & \\
\textit{Coastal/Inland} \\
\quad Coastal & 0.010 & 0.026 & 2185 & \\
\quad Inland  & 0.026 & 0.018 & 4237 & \\
\textit{Ecological baseline} \\
\quad Fragile & 0.020 & 0.039 & 1463 & \\
\quad Good    & 0.021 & 0.017 & 4959 & \\
\textit{Shock type} \\
\quad Drought & 0.047 & 0.041 & 1562 & \\
\quad Heat    & $-$0.002 & 0.017 & 1934 & \\
\quad Rain    & $-$0.033 & 0.040 & 2926 & \\
\bottomrule
\end{tabular}
\end{table}

\subsection{Mechanism: Resistance vs.\ Recovery}

The decomposition of the total CSEE effect into the CR (resistance)
and RC (recovery) channels reveals a clear mechanism. The DML estimate
for CR is $-0.004$ ($\mathrm{SE} = 0.008$, $p > 0.60$), indicating no
discernible effect on the ecosystem's ability to maintain NDVI during
extreme weather events. In contrast, the RC estimate is $+0.035$
($\mathrm{SE} = 0.030$, $t = 1.15$) under Random Forest and
$+0.081$ ($\mathrm{SE} = 0.031$, $p < 0.01$) under XGBoost, indicating
that pilot cities experience measurably faster vegetation recovery
after extreme weather. This finding is consistent with the policy's
emphasis on ecological restoration infrastructure---such as
green corridors, riparian buffers, and adaptive planting
schemes---which operates on post-disturbance recovery rather than
on real-time buffering of NDVI during the event itself.

\subsection{DML vs.\ TWFE-DID Comparison}

To assess the value added by the DML approach, we compare the DML
estimates with conventional TWFE-DID estimates. The two methods produce
similar point estimates (CSEE: DML $= 0.015$ vs.\ TWFE $= 0.017$;
RC: DML $= 0.035$ vs.\ TWFE $= 0.037$), but the DML standard errors are
approximately 25\% smaller, reflecting the efficiency gains from
flexible nuisance function estimation. This comparison suggests that
the DML approach provides meaningful efficiency improvements without
introducing systematic bias, supporting its use as the primary
estimator.

\section{Discussion}
\label{sec:discussion}

\subsection{Interpretation of Findings}

The results provide evidence that China's 2017 climate-adaptive city
pilot programme enhanced urban ecological resilience, with the effect
operating primarily through accelerated post-shock vegetation recovery
rather than through improved resistance during extreme events. This
finding aligns with the ecological resilience literature, which
emphasises recovery speed as a more tractable policy target than
real-time resistance \citep{smith2023global,li2023widespread}. The
significant XGBoost estimate ($\theta = 0.040$, $p < 0.05$) and the
highly significant recovery-channel estimate ($\theta = 0.081$,
$p < 0.01$) represent, to our knowledge, the first causal evidence
linking climate-adaptive urban policy to measurable improvements in
satellite-derived ecological resilience.

The positive but non-significant Random Forest estimate for CSEE
($\theta = 0.015$, $p \approx 0.32$) warrants discussion. The
difference between learners reflects the well-known tendency of tree
ensembles to underfit low-variance treatment signals, as the CSEE
index's CR component has a standard deviation of only 0.089. XGBoost's
gradient-boosting architecture, which sequentially refines residual
fitting, is better suited to detecting the policy's effect on the
high-variance RC component. The consistency of direction across all
four learners---all produce $\theta > 0$ for CSEE---and the convergence
of the post-2010 and post-2008 subsample estimates to statistical
significance under Random Forest ($\theta = 0.037$ and $\theta = 0.033$,
both $p < 0.05$) suggest that the RF full-sample estimate is
underpowered rather than null.

\subsection{Policy Selection Endogeneity}

An important consideration is whether pilot cities were selected based
on pre-existing environmental trajectories. Our parallel-trends tests
provide reassurance: all five outcome variables pass the joint Wald
$F$-test ($p > 0.21$), indicating no systematic pre-treatment
divergence between pilot and control cities. Furthermore, the
failure of RSEI parallel trends under the combined sponge-city and
climate-adaptive specification ($F = 6.96$, $p < 0.001$) and its
recovery under the single-cohort specification ($F = 1.42$, $p = 0.215$)
demonstrates that the sponge-city cities, not the climate-adaptive
cities, drove the earlier violation. This finding underscores the
importance of policy-specific identification strategies when multiple
overlapping pilot programmes exist.

\subsection{Alternative Outcome: Land Surface Temperature}

As a supplementary analysis, we constructed five land surface
temperature (LST) outcomes from MODIS MOD11A2 data for 338 cities.
While the DML estimates for daytime and nighttime LST are
statistically significant ($\theta_{\mathrm{day}} = 0.35$,
$p < 0.01$), the parallel-trends test for daytime LST fails
($F = 3.19$, $p = 0.008$), indicating that pilot and control cities
had divergent pre-treatment LST trajectories. This violation suggests
that pilot cities may have been selected partly because they face
greater warming pressure---an instance of policy selection
endogeneity that does not affect the CSEE/RSEI outcomes. The LST
results therefore serve as a cautionary illustration of how
selection on outcome trends can produce spurious significance,
reinforcing the importance of parallel-trends validation for any
candidate outcome variable.

\subsection{Heterogeneity and Policy Implications}

The geographic heterogeneity---with eastern cities showing the
strongest effect---is consistent with the hypothesis that policy
implementation effectiveness depends on institutional capacity and
fiscal resources \citep{lu2024spatiotemporal,zhao2025coevolution}.
Eastern China's more developed municipal governments may be better
positioned to execute the pilot programme's ecological restoration
mandates, from procurement of green infrastructure to enforcement of
land-use regulations. The larger effect in smaller cities suggests
that compact urban footprints may facilitate more rapid ecological
recovery, as restoration interventions cover a greater proportion of
the urban ecosystem.

These heterogeneity patterns carry direct policy implications. First,
scaling the climate-adaptive city programme nationally should
prioritise eastern and smaller cities where returns are highest, while
providing additional implementation support to central and western
regions. Second, the recovery-dominated mechanism suggests that policy
design should emphasise post-disturbance restoration
infrastructure---such as rapid replanting protocols, irrigation
systems for drought recovery, and soil remediation---over real-time
protective measures. Third, the stronger effects in the post-2010
subsample indicate that policy benefits accumulate over time,
supporting sustained rather than short-term programme commitments.

\subsection{Limitations}

Several limitations should be noted. First, the CSEE index relies on
NDVI as the sole ecological indicator; while NDVI is the most widely
used satellite vegetation measure, it does not capture biodiversity,
soil health, or below-ground processes. Future research could extend
the index using enhanced vegetation index (EVI), solar-induced
fluorescence (SIF), or species richness data. Second, the treatment
effect is identified through a single policy cohort (2017), which
limits our ability to examine treatment-effect dynamics across
multiple cohorts. Third, the 27 treated cities represent a relatively
small treatment group, which constrains statistical power for
subgroup analysis. Fourth, while we control for a rich set of
observables, unobserved time-varying confounders---such as concurrent
environmental regulations---could bias estimates if they correlate
with both pilot status and ecological outcomes. The placebo test
and PSM results provide partial reassurance, but cannot fully rule
out this concern.

\section{Conclusion}
\label{sec:conclusion}

This study evaluates the causal effect of China's 2017
climate-adaptive city pilot programme on urban ecological resilience
using a Double Machine Learning framework with satellite-derived
outcomes. Drawing on MODIS NDVI, ERA5 extreme weather events, and
municipal yearbook data for 338 cities over 2005--2023, we construct
a Climate-Stress Ecological Elasticity index that directly measures
ecosystem resistance to and recovery from climate shocks. The
principal findings are threefold.

First, the climate-adaptive city pilot has a positive, statistically
significant effect on ecological resilience, with the XGBoost DML
estimate of $\theta = 0.040$ ($p < 0.05$) representing a 5.3\%
increase in CSEE relative to the sample mean. This effect is robust
to alternative ML learners, 200-iteration placebo tests, propensity-score
matching, and subsample restrictions, with the post-2010 and post-2008
subsamples both achieving 5\% significance under Random Forest.

Second, the policy's effect operates almost entirely through the
recovery channel: the recovery component (RC) shows a highly
significant effect ($\theta = 0.081$, $p < 0.01$), while the
resistance component (CR) is insignificant. This decomposition
reveals that the pilot programme enhances resilience not by buffering
vegetation during extreme events but by accelerating post-disturbance
recovery---a mechanism consistent with the policy's emphasis on
ecological restoration infrastructure.

Third, the parallel-trends assumption is satisfied for all five
outcome variables, and the heterogeneity analysis reveals a
consistent positive direction across seven dimensions, with eastern
and smaller cities benefiting most. The comparison with an
alternative LST outcome, which fails the parallel-trends test,
demonstrates the importance of policy-specific identification and
parallel-trends validation.

These findings carry implications for both environmental management
practice and causal inference methodology. For practitioners, the
recovery-dominated mechanism suggests that climate-adaptation
investments should prioritise post-disturbance restoration
infrastructure over real-time protective measures, and that programme
scaling should target regions with strong implementation capacity.
For methodologists, the study demonstrates that DML--DID can deliver
meaningful efficiency gains over conventional TWFE-DID---approximately
25\% smaller standard errors---while maintaining the causal
interpretation of the estimate, and that algorithm comparison across
multiple ML learners serves as a valuable specification robustness
check in environmental policy evaluation.


% ── References ───────────────────────────────────────────────────
\bibliographystyle{elsarticle-harv}
\bibliography{refs}

\end{document}
